% ============================================================
%  H. Høffding, Personlighedsprincippet i Filosofien
%  Forelæsninger ved Helsingfors Universitet, foråret 1911.
%  Danish translation (from Swedish) by Mogens Blegvad,
%  ed. Carl Henrik Koch.
%  Scientia Danica, Series H, Humanistica, 8, vol. 6.
%  København: Det Kgl. Danske Videnskabernes Selskab, 2013.
%
%  TRANSLATION (English) — Høffding's six lectures, pp. 11–41.
%  Translator: [HH]
%  Source: scanned 2013 edition (text layer); Høffding's own
%  footnotes only; foreign-language quotations rendered in English.
% ============================================================

\documentclass[12pt]{article}
\usepackage{fontspec}
\usepackage{unicode-math}
\setmainfont{Libertinus Serif}
\setmathfont{Libertinus Math}
\usepackage[english]{babel}
\usepackage[protrusion=true]{microtype}
\usepackage[margin=1.4in]{geometry}
\usepackage{setspace}
\usepackage{fancyhdr}
\usepackage[colorlinks=true,linkcolor=black,urlcolor=black,
  pdftitle={The Principle of Personality in Philosophy},
  pdfauthor={Harald Høffding}]{hyperref}
\setlength{\headheight}{15pt}
\pagestyle{fancy}
\fancyhf{}
\fancyhead[LE,RO]{\thepage}
\fancyhead[RE]{\textit{The Principle of Personality in Philosophy}}
\fancyhead[LO]{\textit{\rightmark}}
\setlength{\parindent}{1.5em}
\setlength{\parskip}{0pt}
\onehalfspacing

\renewcommand{\thesection}{\Roman{section}.}

\title{\textbf{The Principle of Personality in Philosophy}\\[0.5em]
  \large Lectures at the University of Helsingfors, Spring 1911}
\author{Harald Høffding}
\date{Translated from the Danish by [HH]\\[0.4em]
  \textit{Personlighedsprincippet i Filosofien}, trans.\ Mogens Blegvad,
  ed.\ Carl Henrik Koch.\\
  \textit{Scientia Danica}, Series H, Humanistica, 8, vol.~6.\\
  Copenhagen: Det Kgl.\ Danske Videnskabernes Selskab, 2013}

\begin{document}
\maketitle
\tableofcontents
\bigskip


% ============================================================
%  I. PERSONALITY AS THE PRESUPPOSITION OF SCIENCE
% ============================================================

\section{Personality as the Presupposition of Science}
\label{sec:1}
% pp. 11--16

\noindent By the principle of personality I understand the thought that the
human being, at every point and in every respect, is to be treated---by others
and by himself---as an end, never merely as a means. Even where the human being
serves---serves other human beings, serves society, serves a cause, an
idea---he is at the same time to be an end himself, in that this serving ought
to lead to the development and ennoblement of his own being, of his own
individual distinctiveness. It is an ideal that is set up in this principle,
which was first formulated with full clarity and definiteness by Immanuel Kant
(1785). In the Nordic countries it was asserted with great emphasis by Erik
Gustaf Geijer in Sweden (1842), from whom the term ``the principle of
personality'' derives. Whereas Kant set up this principle in opposition partly
to the principle of authority, partly to an external morality of utility and
happiness, Geijer set it up against the romantically conservative ``historical
school,'' which laid weight solely on institutions and traditions as they
happened to have developed, and he found more and more that this principle
contained everything for which he had worked---in literature, in politics, and
in religion. In Denmark the principle is expressed, particularly in the
religious domain, at roughly the same time (1846) by F.C.\ Sibbern and Søren
Kierkegaard. --- If we were to follow the history of the principle further back
in time, we would come on the one hand to Christianity and the weight it laid
on the individual soul, on the other to the Socratic insistence on the
importance of self-knowledge and self-examination.

The justification of the principle is very different among its various
representatives. I shall return to this in connection with the inquiry into
whether the principle admits of any binding proof at all. Here I will attempt to
set forth all that it contains. It is first and foremost set up as an ethical
and religious principle; but I believe that it also contains points of view of
great importance for a correct conception of the relation of science to life.
It has intellectual, not merely ethical and religious, significance, and I shall
examine it from all these three sides. What Geijer said of his age---that it
``is an age that suffers from its own thoughts,'' in that ``the inner world
[\ldots] gives way in its very foundations, in its faith, its conviction, its
knowledge,'' and that ``the fundamental and mother-thought that here strives to
come forth and make its way, and under whose birth-pangs the world
trembles''\footnote{E.G.\ Geijer, \emph{Samlade skrifter}, sec.~I, vol.~6,
  Stockholm: P.A.\ Norstedt \& Söner, 1853, p.~287.} is precisely the principle
of personality---this surely holds true to this very day. The crisis is not
over; it has merely assumed somewhat different forms than in Geijer's time.
Geijer's words aptly express that here is a thought which has involuntarily
worked its way forward, but which can come into its full right only by becoming
an object of clear consciousness.

When we now begin with the intellectual significance of the principle of
personality, we meet the objection that in science it is precisely a matter of
disregarding our personality---its striving, its pleasure and pain, its hope and
its fear. Ought we not precisely to divest ourselves of our personality if we
would be researchers, if we would find the great laws of existence that hold for
each and all? In science we are led into a common world. But every personality
has its own little world, which is different from everyone else's, and which,
like the dream-world, stands in opposition to the waking reality common to all.
Can the personality here demand more than to be allowed to live in its world of
dreams, so long as science does not draw this world under the strict law of its
research?

I enter for a moment into this analogy between dream and reality on the one hand,
personality and science on the other. Psychology shows that the same laws hold
for consciousness in dreams and in the waking state. In both there is revealed a
need for, and a capacity for, comprehending and ordering all experiences, all
sensations, all feelings, and all thoughts. In dreams this happens in a more
abrupt and apparently irregular manner---now ingeniously, now stupidly. In the
waking state it happens more according to definite theoretical and practical
points of view (although here too both stupidity and ingenuity can occur).
Analogous is the relation between personality and science. It is a kindred need
and a kindred capacity that is at work in both. We speak of personality only
where there appears a striving for coherence in the life of the soul, a striving
to gather all psychic elements and to direct them in definite directions.
Continuity and concentration are the special marks of personality. Science is
merely a definite, specially developed form of this distinctiveness. In so far
as there exists an opposition between personality and science, it must be
established that it lies within the personality---that it concerns the relation
between the personality as a whole and one of its special expressions of life.

It is precisely one of the chief tasks of philosophy to investigate how and why
scientific consciousness develops as a special form of personal life.
Philosophical reflection seeks to answer the question of how the nature of the
human mind expresses itself both in the work that leads to a great, common
world-picture, and in the work that perhaps acquires significance merely by
bringing harmony and order into the individual's little world. This little world
need not be a dream-world; here the analogy mentioned above cannot be applied
without further ado---at any rate there is here a special problem to which we
shall later return.

It is the critical philosophy that has, with the greatest energy, maintained
that the work of thought in science cannot be completed, and cannot be
understood in its full significance, without taking into account the nature and
conditions of conscious beings. Every objective truth has a subjective
foundation, even though one cannot infer this truth from the merely subjective
foundation.

The historical origin of science is to be found wherever the human being has had
the experience that he can attain his ends only along quite definite paths. He
would rather attain them immediately. Only when compelled by experience does he
take roundabout ways. And he then stands before a connection to which he must
bow. Here occurs his first encounter with the necessity that science discloses
at one point after another. Science demands resignation. The human being here
thinks in order to live. Thought, and thereby the work of personality, is here
merely a means. But it is a means that can itself become an end by becoming an
object of immediate interest. It becomes an independent task to come to know the
great connection within which personal life is lived, to find an order of things
more comprehensive than the individual's own little world. And then the human
being can live in order to think. He grows with the more comprehensive ends he
thus sets himself. The personality becomes richer the larger a circle of
experiences it embraces, provided only that it can gather them into a harmonious
unity, or at least divine and seek such a unity. It is easy enough to
concentrate when there is only a poor and limited content to comprehend, when no
great oppositions, no wealth of experiences, lay claim to the psychic energy.
The world of the personality must be expanded for this energy rightly to show
itself.

Science and its history become intelligible only when there exists an inner
connection between scientific work and the whole life and activity of the
personality. Scientific work presupposes an energy and an enthusiasm, an
\emph{amor intellectualis}, which is not possible without an inner connection
between science and personality. And such a connection presupposes, in turn,
that it is no alien power to which the personality bows when, through research,
it is led to recognize a great connection that points far beyond its own little
world. The longing to find truth along the path that science indicates is merely
a special form of the longing for agreement with oneself amid all manifold and
changing experiences---the longing in which the personality's own being
consists. And even where the personality keenly senses the opposition between
its own little world and the unsurveyable world---reaching beyond all human
action and ends---for which science opens the eyes, it belongs precisely to the
nobility of the personality, as well as to its duty, to find by its own power,
by the exertion of its own thought's energy, the great laws that condition its
existence and its worth. The personality is purified by seeing its existence
conditioned and limited by the great connection of things.

Scientific work does not lay claim in equal degree to all sides of the
personality. In his research the human being must proceed from definite points
of view, follow definite paths, develop his life of thought according to certain
definite methods, which, to be sure, are special forms of personal life in
general, but which in the course of mankind's long labor of thought have proved
their indispensability. Science works with certain fundamental concepts or
fundamental predicates, which from of old are called categories, and which
together constitute an armament that adapts itself both to the personality and
to the resistance that the work of thought is to overcome.\footnote{Cf.\ H.\
  Høffding, \emph{Den menneskelige Tanke, dens Former og dens Opgaver} [The
  Human Mind: Its Forms and Its Tasks], Copenhagen and Kristiania: Gyldendalske
  Boghandel, Nordisk Forlag, 1910, ch.~3: ``Tankens Former (Kategorier)'' [The
  Forms of Thought (Categories)], pp.~135--245.} At times one has called the
personality, in so far as it appears in this armament, ``the pure I,'' in
opposition to ``the empirical I,'' which shows us the personality in its whole
individual reality, with all its experiences, representations, wishes, and
drives. The pure I is the I of the solitude of thought, of the study and the
laboratory, the I that works strictly with the ``categories.'' But it is merely
a particular formation---necessitated by the division of labor---of the same I
that manifests itself in representations, wishes, and drives. There is no
absolute difference between them, and their mutual relation can appear in
countless nuances in different individuals. There are researchers who are
transformed into wholly other people when they step out of the solitude of
thought into ordinary life. But this is an imperfection that points to an
unresolved disharmony in personal life itself.

On this point stands the dispute of our days between criticism and pragmatism.
The first is (already in its founder Immanuel Kant) inclined to regard ``the
pure I,'' the sum of the forms of investigating thought (the categories), as
given once and for all, independent of all experience and history (including
science's own history). The latter is (both in the form Ernst Mach and in the
form William James have given it) inclined to regard the forms and
presuppositions necessary for scientific work as merely conditioned by external
tasks and practical interests, without their standing in any inner connection
with human nature itself. In my view this contradiction falls away when one more
closely examines the nature and development of the individual forms of thought.
In my book \emph{The Human Mind} I have worked in this direction and have
thereby been strengthened in the conviction that the human being in science
neither stands merely as a means for the application of certain pure forms nor
for the solution of externally given tasks.

When science thus, in its origin, points back to the human personality and
appears as a special form of the work that every personality has to perform in
order to assert and develop itself, its personal character becomes visible from
another and more individual side at the conclusion of scientific work. An
absolute conclusion, a conclusion once and for all, is not possible, because
life and experience continually go further. But each individual researcher shall
and must stop at a definite point, where inner and outer causes prevent him from
going further. And then he must of course seek an overview, a summation, a
rounding-off. He must also attempt a glimpse into the unexplored, if not
unexplorable, regions, just as one casts light into the darkness with an electric
searchlight. Here it is not merely ``the pure I'' that works; the other sides of
the personality must more or less take part. One sees this, for example, in the
natural scientist's attempt to set his results in relation to his whole view of
the world. And one sees it when the historical researcher passes from critical
scrutiny of the sources to the writing of history; this requires a rounding-off
in which the researcher's personality asserts itself more than during the
critical work. One sees it still more in philosophy, when it passes from
criticism to system, from analysis to synthesis.

Philosophical systems ought never to have wished to be more than an individual
rounding-off. A rounding-off can be justified even though a conclusion is
impossible. Goethe characterized his research with the principle: ``Nie
geschlossen, oft gerundet!'' [Never closed, often rounded!] And in the
rounding-off the individual character ought to be expressly acknowledged. The
truly significant systems lose nothing thereby, but acquire through it a special
interest. Even the system that is often regarded as the most objective and
impersonal of all---Spinoza's---proves on closer examination to be permeated by
its founder's personality; it constitutes the band that unites the many lines of
thought into a unity. Through whole trains of thought one feels the master's
undaunted consistency, his great resignation, and his heartfelt absorption in
the work of thought, which for him was a religion. It is not least this that
makes many return again and again, and ever with new profit, to this artwork of
thought.

The interest that the history of philosophy presents rests for a large part on
this personal element. The history of philosophy can be regarded as a great
dialogue, whose individual lines stem from researcher-personalities, each of
whom, on the basis of the science and culture of his age, has developed a type
of thought, a more or less essential view of existence. For philosophy, its
history has, precisely because of the principle of personality, greater
significance than for other sciences. It stands, in this as in other respects,
between science and art. And for many it is precisely herein that its great power
of attraction lies.


% ============================================================
%  II. PERSONALITY AS AN OBJECT OF SCIENCE
% ============================================================

\section{Personality as an Object of Science}
\label{sec:2}
% pp. 17--21

\noindent For a correct conception of science one must see in it a fruit of
personal work. Thereby both the fear of its results and pride over them are
excluded.

But the personality is not merely the subject of science. It is at the same time
itself one of the objects of science. Science turns toward it as toward a given
subject-matter, whose nature and origin are to be understood and explained. The
personal factor that co-operates in all science can and ought itself to be made
an object of science. The world is understood through the human being
(\emph{ex analogia hominis}); knowledge of the world is conditioned by the
capacities and forms of human thought. But the human being is himself a part of
the world, which can be understood only in its connection with the rest of the
world. The human being is thus understood through the world
(\emph{ex analogia universi}).

As Pascal has said: ``In space the universe comprehends (\emph{comprend}) me and
engulfs me as a point; but in thought I comprehend (\emph{comprends}) the
universe.''

The two modes of consideration do not conflict with one another. They are two
points of view, both of which must be adopted, and which stand in inner
reciprocal action. Human thought, which works to solve the riddles of existence,
has itself developed within existence. The presuppositions that condition and
limit our conception of the world have themselves their history. It is not least
the necessity of this double point of view that makes an absolute conclusion of
knowledge impossible. We cannot apply both points of view at once. We cannot at
one and the same time stand on our head and on our feet; but we can (perhaps)
first do the one and then the other.

Against making the personality an object of science, however, one of the finest
thinkers in present-day Germany, Heinrich Rickert, has raised an interesting
objection in his work \emph{Die Grenzen der naturwissenschaftlichen
Begriffsbildung} (1901--02). Rickert defends a sharp opposition between natural
science and history (the science of mind); the first concerns the universal and
general, that which constantly repeats itself---the latter has for its object
something that happens only once (\emph{das Einmalige}), individual events and
individual figures. The purely natural-scientific method aims at finding general
concepts and laws under which the individual phenomena can be subsumed as
examples. The historical sciences, by contrast, stop at individual totalities.
The natural-scientific concepts are general concepts, the historical concepts
are individual concepts.\footnote{``Von den Individuen gibt es keine
  Naturwissenschaft, weil sich nichts Besonderes mehr ihnen unterordnen lässt,
  in Vergleich zu dem sie noch als ein allgemeiner Begriff oder als eine `Natur'
  aufzufassen wären.'' [Of individuals there is no natural science, because
  nothing more particular can be subsumed under them, in comparison with which
  they could still be conceived as a general concept or as a ``nature.''] H.\
  Rickert, \emph{Die Grenzen der naturwissenschaftlichen Begriffsbildung. Eine
  logische Einleitung in die historischen Wissenschaften}, Leipzig: J.C.B.\
  Mohr, 1901--02, p.~295\,ff.}

Against this conception I set the assertion that in nature as well as in the
world of mind, general concepts and laws are in the last resort merely means for
understanding the individual and historical (\emph{das Einmalige}). It is the
pride of science that it can form concepts and laws that hold for countless
individuals and events. The capacity for this was already, for Plato, the
typical figure of Greek science, a divine gift that ever anew filled his soul
with admiration. Indeed, for Plato even the gods were subordinate to the eternal
ideas and rejoiced to look up to them. But no actually given phenomenon, neither
in nature nor in the world of mind, is exhaustively explained by being regarded
as an example of a general concept or of a general law. To understand natural
phenomena an interplay of laws is required. The laws are the instruments of
thought that, through their co-operation, make the arising and development of a
phenomenon intelligible. The solar system, the Earth, the organism, the
personality, society---these are examples of individual totalities, whose history
it is the task of science to write by means of the concepts and laws it is able
to find. The history of science shows us clearly the increasing significance of
the concept of development. In the end we are confronted with the task of writing
the biography of `the World'---in so far as `the World' can take shape as a
rounded whole.

It is by no means the case that it is only history which exhibits the singular
phenomenon (\emph{das Einmalige}). Neither does absolute repetition occur in
nature. No event ever recurs in exactly the same way. The immutability of nature
is only apparent. It is a great world-drama (or perhaps rather an epic) that
natural science and history together are to interpret. In it no scene, no
situation, no line occurs more than once, if one carefully notes all the
particulars, all the nuances and accents.

It is a great and lofty ideal that is hereby set before science. But it is
necessary for rightly understanding the significance of science. Natural science
approaches it by means of the great hypotheses that are raised on the basis of
the laws that have been found. The laws always hold for simple and elementary
relations; then it becomes a problem how these relations assert themselves within
the wholes of various kinds and various degrees that experience shows. Whereas
natural science goes from the elements to the totalities, the science of mind
stands more directly before totalities (events, personalities, societies), and
such as do not allow themselves to be resolved into elements as easily as the
material totalities that natural science treats. This, however, conditions only a
difference of degree, not a difference of kind.

It is a long-familiar matter that the more comprehensive our concepts are, the
poorer they are in content, and conversely. But it is not the richness of content
that is to be subordinated to the comprehensive abstractions; it is these that
are to be applied in order to understand and appreciate the former. The most
universally valid sciences, logic and mathematics, are applied (apart from their
own independent interest) as means and presuppositions for the more special
sciences. The difficulties increase and the riddles become more numerous the
more we approach the individual totalities. But these still indicate the ultimate
objects and tasks of research.

When it is said that we cannot form proper concepts of the individual
(\emph{das Einmalige}), because there is nothing that can be subsumed under it,
the answer must be that to the concept we form of an individual phenomenon, an
event (e.g., the French Revolution) or a historical personality (e.g., Luther),
belong all the various sides, properties, situations, and degrees of development
that can contribute to its complete characterization. The ``historical''
concepts too, not merely the natural-scientific ones, comprise a manifold. Two
kinds of individual concept must be distinguished. If I think of a historical
personality in a definite situation (e.g., Luther at Worms), I can form of him a
concrete individual concept; it acquires an intuitive character and bears the
stamp of the memory-image or the fantasy-image. But no historical personality is
exhausted in a single situation (any more than in a single property or a definite
phase of development). My complete concept of it must correspond to all the
situations in which it appears, and this may then be called a typical individual
concept. The great difficulty with individual concepts is precisely to determine
the relation between the concrete and the typical. Is it the young Luther, he who
burned the papal bull, or is it the old Luther, the organizer and dogmatist, that
is to be decisive for our final concept of Luther? Which of the phases of the
French Revolution are to determine for us its (typical) concept?

The solution surely lies in this: that the typical concept of the individual
indicates the law according to which the transition takes place from the one
stage and the one situation to the other. It is the law of change that gives
expression to the essence of an individuality. If this law can be discovered, the
nature of the individual must be rediscoverable in each of its stages, and one
must be able to trace in the preceding stage the one that follows.\footnote{Leibniz,
  the real founder of the individual concept, already expressed this thought. It
  is ``la loi de l'ordre'' [the law of order], ``la loi du changement'' [the law
  of change], that constitutes the individuality of each being. \emph{Lettre à
  Basnage de Beauval}, 1698 (G.W.\ Leibniz, \emph{Opera philosophica}, ed.\
  J.E.\ Erdmann, Berlin: Eichler, 1840, p.~151).} Every psychological and
historical presentation works to find such a law. Every good biography at least
lets us divine it.

Art can show us only the concrete picture of individual situations and actions in
which the personality appears, not the whole series of pictures that would be
necessary to form a typical concept. Therefore contradictory interpretations of
artistically presented figures can arise, e.g., of Hamlet. And yet art is able to
carry through the individualization better than nature and history as a rule do.

However difficult it may be to form typical individual concepts, these
nevertheless stand in the end as the highest task of thought. In order to be able
to find general laws, science must often disregard the fact that everything
experience shows us has an individual character, appears as a whole, even if
perhaps a simple and elementary one. There are many degrees of
individuality---from a cloud in the sky or a heap of sand on the shore to a human
personality or a human society. All too often one meets the notion that nature
``begins'' with a chaos of elements and forces without quality and definite
direction, and that all totalities and individualities are products of these. In
our textbooks we begin with the simple and then pass to the composite, because
understanding is most easily attained in that way. But what we know about the
elements, about the simple, is in science always won through analysis of the
total connection, by abstraction from individual phenomena that can vary to
infinity.\footnote{Cf., with respect to the question treated in this section, my
  treatment of ``Totality and Development'' as special categories in \emph{Den
  menneskelige Tanke}, Copenhagen: Gyldendalske Boghandel, Nordisk Forlag, 1910,
  pp.~218--37.}

The consideration here set forth acquires special significance when the
personality becomes an object of science. For the personality is not merely the
example, nearest to us, of an individual totality. It is at the same time the
most pronounced example of such, in that it displays the most intimate
co-operation between elements and forces that lies within our experience.

We now turn to psychology, the empirical science of the personality, in order to
see what it can teach us about the essence of the personality in itself.


% ============================================================
%  III. THE PRINCIPLE OF PERSONALITY IN PSYCHOLOGY
% ============================================================

\section{The Principle of Personality in Psychology}
\label{sec:3}
% pp. 22--26

\noindent In all empirical sciences the decisive definitions come at the end,
not at the beginning. Psychology slowly works its way toward a concept of
personality. All its various methods---observation, description, comparison,
experiment---lead directly or indirectly to this goal. When we already now wish
to attempt to set up---and in what precedes have presupposed---a concept of
personality, this can be done only by way of hypothesis.

As a characteristic trait of personal beings must first be brought forward the
uniformity in all states and undertakings. We call a human being a personality
in all the more pronounced a degree, the more its thought, feelings, and
endeavors are directed toward one goal---the more concentration there is in its
spiritual life, the greater the continuity there is in its striving. Despite all
changes and changing situations, there must appear in it a common law, a stamp
of unity, or at least a coherence. This shows itself in the most exalted manner
when the human being is subjected to great fluctuations, finds itself in strong
conflicts, and yet, even in the dark hours, does not deny what, at the high
points of life, stood before it as true and right.

The personality is maintained---like all life of the soul---only by an
uninterrupted spiritual work, in which the psychic energy shows itself. It is a
constant struggle to hold together and harmonize a given manifold of sensations,
thoughts, feelings, drives, and interests, which go in different directions or
are indifferent to one another. The greater this manifold and disparity, the
more spiritual work is required for the life of the soul to subsist. And the more
intimate the coherence and concentration, the greater the psychic energy that
must be applied to this work.

It was this concept that I took as my basis when, about thirty years ago, I
worked out my exposition of psychology. It had developed in me through criticism
of the older English psychology (the so-called association psychology), which
sought to explain the life of the soul---and thereby the personality---as the
product of a mechanical joining-together of psychic elements (independent
sensations, representations, and so on), which stood in roughly the same relation
to the personality as a heap of stones to the house that is built. The task now
became to show that there is never, so long as the life of the soul subsists, a
complete lack of coherence and unity in it, although this coherence can assume
very elementary forms. I supported myself here especially on the doctrine of
mental illnesses, which proves that mental illness consists specifically in a
dissolution of the coherence within the life of the soul, in a weakening of the
psychic energy. Later my conception was confirmed by Pierre Janet's work,
especially by his description of ``psychasthenia,'' a spiritual weakness
characterized by the fact that no psychic content, no psychic elements, have
fallen away, but that there is a lack of energy for cohesion and concentration,
on account of which patients suffering from this illness are unable to make
decisions, to take resolutions. There is lacking what Janet aptly calls
``psychic tension'' (\emph{tension psychologique}), which corresponds to what I
call synthesis or psychic energy.

In the child the psychic energy is not yet so developed that a strong and lasting
concentration can become possible. Spiritual development consists not merely in a
manifold content's being taken up through experience and tradition, but first and
foremost in the energy that brings about an inner coherence. The dream-state is
characterized especially by the lack of firm coherence and consistency in the
content of consciousness. Like the child's consciousness, it has a sporadic
character. In it one representation easily passes over into another; the
situations change without this, as a rule, arousing astonishment and without an
explanation of this alternation being sought. In the waking consciousness
definite directions are held fast and pursued, and the changes are motivated, or,
if motivation seems to be lacking, become riddles. Even if there is here only a
difference of degree---since there are also waking dreams---this difference
nevertheless clearly shows the significance of the psychic energy, especially
when we have reason to assume that the dream-consciousness denotes a state
between the conscious and the unconscious.

It is, however, not merely from weakness or lack of development that the
individual elements of the life of the soul---sensations, thoughts, feelings, and
so on---can appear independently and without closer coherence. We can also, by
attention and analysis, bring forward and isolate individual elements and, as far
as possible, consider them each by itself. We must distinguish in order to be
able to apprehend distinctly. But such distinctions are more or less artificial.
They are reflections, the work of after-thought, and must not be confused with
what immediate experience offers. The more we are able to immerse ourselves in
our purely immediate and involuntary experiences, the more clearly will our life
of the soul show itself to flow as a continuous stream, in which no element can
be isolated from the others without changing its character. This continuous
character of the immediate psychic experiences has in the most recent time been
demonstrated in a brilliant and emphatic manner by William James and Henri
Bergson.

Finally, experimental psychology too points in the same direction. Its real
founder, Fechner, proved that the individual sense-sensations do not arise
independently of one another, but are interwoven with one another, in that after
a very strong stimulation of the sense-organs a stronger stimulation than usual
is required for a sense-sensation of a certain strength to arise. The so-called
contrast-effects too bear witness to the great coherence, even in our elementary
life of the soul. The green color appears more pronounced (more saturated) after
the red, which is its complementary color, than after other colors. The absolute
independence of the sensations is merely apparent.

And if we consider the matter from the physiological side, we are led in the same
direction. However the relation between soul and body may be, it nevertheless
stands fast that the life of the soul is exceedingly intimately bound to the
nervous system. But the significance of the latter lies in this: that it is a
great organ of concentration, through which all parts of the body stand in close
connection and reciprocal action, so that thereby a unified and concentrated
performance of the organism as a whole is possible. And the higher and the more
energetically developed a life of the soul we have reason to assume in a living
being, the greater the role the nervous system plays, and the more centralized it
is. In the outer, sensible appearance too the distinctive character of the life
of the soul, and thereby of the personality, manifests itself.

It is only in those beings whose life of the soul has attained a higher degree of
unity and clarity that we speak of personality. Personal life can appear in
inconceivably many different degrees; for a large part the difference between the
various individuals rests on such a difference of degree in the intensity of
concentration. Only in human beings does one speak of personality. The life of
the soul of animals must, to judge by all experience, be assumed to be partly so
simple, partly so changing and incoherent, that we must picture it most nearly in
analogy with our dream-life.

In what precedes we have defined the concept of personality solely from its
formal side. We have spoken of energy and concentration, of unity and coherence.
But the content too, which is gathered into unity, has its significance. Two
personal beings with the same capacity for concentration can nevertheless be very
different on account of the different content they contain, and around which they
gather themselves, as well as the different direction in which their energy is
released.

In the real respect, the personality acquires its stamp from the fact that it
embraces a central content, a ruling interest and striving, which conditions the
direction in which the concentration takes place and the psychic energy is
applied. Each individual experience, thought, mood, decision acquires within the
personality its definite place by its relation to a fundamental value that
constitutes the real distinctive mark of an individual personality, just as the
degree of the energy of concentration constitutes its formal distinctive mark.
The little world that each personality embraces has at every time and at every
stage of development a center, in relation to which all other elements in the
personal life are determined with respect to value and significance. To know
oneself is to have consciousness of this central interest, this fundamental
value, which constitutes what one might call the real I. If there were no such
constantly active elements in us, we would be strangers to ourselves. There can
be constant elements of a purely peripheral kind in us; but the deeper
self-knowledge presupposes central elements, and it therefore makes a comical
impression when Holberg's Jeppe recognizes himself by his hollow tooth. Such a
peripheral element cannot determine the place of all other elements in the inner
world-order.\footnote{On the difference between the formal and the real I, cf.\
  my \emph{Psykologi i Omrids paa Grundlag af Erfaringen} [Outlines of
  Psychology on the Basis of Experience], Copenhagen: Philipsen, 1882 (and
  later), V~B, 5.}

In personal life a crisis arises if the central elements are changed. This can
happen through sudden catastrophes or through long periods of doubt and unrest;
but it can also happen slowly and imperceptibly, so that only afterward do we
discover that an essential change has taken place with us and in us. Something can
work its way forward in us that quite involuntarily determines our striving, but
of which we become conscious only later. An example that lies near at hand here
is Geijer's relation to the thought to which he gave the name the principle of
personality. It arose in him at a late point in his life, in his fifty-fifth
year, but he then discovered that he had long followed it without himself knowing
it. In a letter that he wrote on his last birthday (1847), he expresses himself
about it in the following manner:

\begin{quote}
There is a fundamental thought in my whole life, which steps forth ever more
clearly, and which is not my own work; for it has led me and carried me right up
to this day. Were I to give it a name, it is none other than the much-disputed
principle of personality, whose prophet I have become almost against my will. I
mean by this that I have held off from myself and shaken off everything that
could hinder and hamper the development of my own innermost being. This has for
the most part happened without my own co-operation. But it has now come over me,
and it becomes clear to me that what I have fought for in literature, in
politics, and in religion is properly this my innermost personality and its
right, in the best sense.\footnote{E.G.\ Geijer, \emph{Samlade Skrifter},
  sec.~I, vol.~1, Stockholm: P.A.\ Norstedt \& Söner, 1849, p.~XXXVIII.}
\end{quote}

Thus the very developmental history of the principle of personality presents an
example of a chief point in the nature and conditions of personal life. We are
led to see that the personality cannot be said to be exhaustively determined by
what, at a given moment, appears with clear consciousness and full understanding.
Just as one assumes that in organic life apparently sudden changes (mutations)
can occur, through which new types can arise, so too in the spiritual domain,
through a quiet increase, changes in the lines of development can be prepared
which to the outer observer may seem inexplicable, although they have their cause
in what has taken place in the innermost core of the life of the soul. We have
here a real example of the difference between concrete and typical individual
representations that was set forth in the preceding section. The typical concept
of a personality can be reached only by following it through its various stages of
development. And it is often the later stages that draw attention to the content
and significance of the earlier ones. Descriptive psychology has here a great
field before it. And it must be hoped that here, when the sense for personal
life---its forms and phases---has been developed still more than hitherto, a
significant work will be able to be carried out. Hitherto, dogmatic and
doctrinaire conceptions of various kinds have hindered an uninterrupted work in
this direction. Descriptive psychology approaches the boundary between science
and art; but by clarity in conception and by the application of the comparative
method it can conquer new terrain for science.


% ============================================================
%  IV. THE PRINCIPLE OF PERSONALITY IN ETHICS
% ============================================================

\section{The Principle of Personality in Ethics}
\label{sec:4}
% pp. 27--31

\noindent From the intellectual significance of the principle of personality I
now pass to its practical significance, which is the original one.

Every principle sets up a demand to our thought or to our will. The principle of
personality makes demands in both directions. I have attempted to show what
demands it makes on the life of thought. The practical demand that it makes
becomes, as I shall now attempt to show, of great significance for ethics and the
philosophy of religion.

In Kant and Geijer, in Kierkegaard and Sibbern alike, the principle of
personality is set up in a critical and polemical direction against the dominion
of society and tradition over the individual personalities. It demands their free
and independent development. But it nevertheless contains within itself the
possibility of letting both the social and the individualistic conception of human
life and human action come into their right.

From a purely social point of view it is emphasized that the individual always
arises within a society. An isolated individual does not exist. The individual
conscience itself, the innermost in the human being, is a product of social life,
a social inheritance. The individual virtues are socially conditioned. And on the
other hand the individual, with its longing and its striving, is vanishing in
comparison with the life of society, which continues through countless
generations. It is merely a drop in the ocean. Its happiness, its knowledge, its
conscience acquire vanishingly slight significance in contrast to the great and
constant course of development of the race and of society. And the only measure
of the individual's significance lies in the contribution it is able to make to
this course of development.

The individualist answers, for his part: Where should the significance of society
be sought if not in the security and support it gives to the development of the
individual personalities? The worth of a society is measured by the degree to
which it fulfills this task. This is the only possible measure, already for the
good reason that society consists of individuals and exists only in and through
them.

If the social conception makes the individual a means for society, then the
individualistic conception makes society a means for the individual.

This opposition does not come forward so long as life is essentially lived under
the involuntary guidance of instinct and tradition. But every cultural
development comes to a point where the individual no longer feels itself in
solidarity with society, but where society and its demands stand as mere
tradition, mere authority, perhaps as purely external power. It is such an age
that, in Geijer's words, suffers from its own thoughts. There have then---precisely
as a consequence of social development---developed independent personalities; but
society does not yet fully dare to acknowledge their independence, and they
cannot acknowledge the worth of society. It is really only in such times that
philosophy arises. Philosophy therefore acquires, in its origin, an
individualistic character. Greek philosophy, the thought of the Renaissance, and
the philosophy of the eighteenth century are significant examples of this. The
philosophy of the nineteenth century denotes---both in its romantic and in its
positivistic main form---a reaction against this, just as the great labor
associations too, whose coming into being belongs among this century's greatest
events, have an anti-individualistic character. But the problem arises again and
again, and we have entered the new century with it.

With respect to this problem a sharp distinction must be made between the
historical and the philosophical mode of consideration. It is a fact that the
individual is developed within society. Its very longing and its striving, its
ethical conceptions and the criticism it exercises toward tradition, cannot be
understood if one starts merely from the individual, thought of as isolated. The
individual is to that extent a product of social causes. And society extends far
beyond the individual, both into the future and into the past. But does there
follow from this, with necessity, an obligation for the individual to bow under
society's laws and to work for its welfare? The individual can, after all, regard
the social development, whose fruit it itself is, as a mere fact, a result, which
it---if it has the spiritual and mental power to do so---now applies in the
service of its own individual ends. It serves no purpose to speak of a social debt
that the individual is to pay off by working for society's continued development,
for the question is precisely why the individual should acknowledge this debt.
Why should the individual not regard itself as sovereign?

What individualism in its extreme forms asserts is really this: that the
personality is to be an end without being a means. Against the principle of
personality, by contrast, it does not conflict that the personality is a
means---that it serves a cause, an idea, serves society---provided only that this
serving is a means to a higher development of the individual itself. The human
being grows, as Schiller has said, with his greater ends. But these greater ends
it can find only in social life. Just as thought grows only by seeking truth---that
is, by finding as great a connection as possible among as many experiences as
possible---so the life of feeling and will grows only through the experiences that
are made and the work that is performed by participation in the life of society
and of the race. But the condition for this is that society's forms and demands be
modified, so that the social work---the paying off of ``the social debt''---can
serve to bring to development powers and possibilities that would otherwise
slumber unused in the individual. The personality can make itself a means for
society only when society also makes itself a means for the personality. The
possibility of this is the kernel of the so-called social question. In all of
history we find great groups of human beings who have essentially been treated as
means, without society's regarding itself as having duties toward them. All
justified social reforms aim at abolishing this dualism of end and means. Society
must, if it is to continue to subsist, acknowledge the principle of personality,
even though this was originally enunciated from individualistic standpoints.

Some forms of individualism, however, overlook this view. Max Stirner proclaimed
(in his book \emph{Der Einzige und sein Eigenthum}, 1845) a philosophy of the I
as abstract as it was absolute, according to which there is, neither within nor
outside us, anything whose claim we are to acknowledge. So long as you believe in
the truth, he even says, you do not believe in yourself alone and are only a
servant. Every right is a privilege. I acknowledge no right, and I do not care to
have right, if only I have might!---Here society is not merely made into a means;
it disappears entirely. In contrast to this, an individualism or ``anarchism''
such as Kropotkin's rests on the experience of the significance that involuntary
association has for all living beings, and he has depicted it in an interesting
manner in his book \emph{Mutual Aid} (1902). Kropotkin merely opposes
centralization, the state's intervention in the spontaneous social formations in
smaller circles. In neighborly relations life orders itself as if of its own
accord. What Kropotkin does not have in view is that these involuntary
associations---like the earlier guilds and fellowships---can lead to stagnation in
oppression and injustice, and that therefore an intervention can be needful. But
he is right in this, that the state ought not to intervene unnecessarily in the
smaller circles in which life stirs of its own accord---for the same reason that
society ought not to intervene in the free development of the individual
personality, unless this becomes necessary out of regard for the free development
of other personalities. It is in the interior of the individual that the new
powers which can carry social life forward have their origin. The growth of the
future sprouts there. It is thus for its own sake that society ought to
acknowledge the principle of personality. Its real power shows itself in the
strength and independence with which the individuals can live within its frame.

One often meets, among modern sociologists, a certain mystical conception of the
concept of society, in that they assume a social origin for ideas and feelings.
If by this is merely meant that ideas and feelings are developed in the individual
through reciprocal action with others, it is completely correct. But when
``society'' is invoked as a cause, it is mysticism. Society as such can produce
right, morality, and religion as little as it can produce poetry and science. The
``folk-songs'' were not composed by ``the folk,'' but by individual singers whose
names are forgotten because their work at once became everyone's property and, by
passing through many hands, underwent changes that often made it into something
other than it originally was. At earlier stages of culture individual initiative
is easily forgotten. And it is precisely such stages of culture that the
sociologists prefer to occupy themselves with. At later stages the founding
personalities play a greater role and are not so easily forgotten. An appeal to
``society'' as the explanation of a phenomenon is in reality an \emph{asylum
ignorantiae}.\footnote{It is in this connection of interest that E.G.\ Geijer's
  defection from ``the historical school,'' which was the romantic precursor of
  modern sociology, was precisely brought about by his historical studies (E.G.\
  Geijer, \emph{Samlade skrifter}, sec.~I, vol.~8, Stockholm: P.A.\ Norstedt \&
  Söner, 1855, p.~520). That is to say, he recognized that it was not an
  explanation of a phenomenon to appeal to ``society'' or ``history'' \emph{en
  bloc}. And he recognized ``that history cannot give, but can indeed modify, the
  principle of action, and that the latter must, on the contrary, be the
  contribution of each age itself.''}

With this, the justification of setting society and individual in opposition to
one another is not to be denied. The concept of society contains the fact that, by
means of the life and work of changing generations, a capital is gathered, a
potential energy, which is deposited in institutions and traditions, indeed in
language itself. Here is something that points beyond the individual---back into
the past and out toward the future---something that belongs to no single
individual whatever. They are values that are bound, and that can be released only
through the work of individuals. A social order is valued not merely according to
the capital it thus has at its disposal, but also according to the manner in which
it is able to release it. There can be dead values that acquire the semblance of
life through social institutions. Whether the values are living or dead can be
decided only by whether or not they benefit the development of free personalities.
For the free personalities are both the centers in which the values of life reveal
themselves and the centers from which the work of finding and producing values
proceeds.

According to this whole view, society and individual stand at the same time as
end and means for one another.

The highest ethical thought is this: that the individual personality is to develop
in such a way that precisely by its independent and distinctive development it, in
the best manner, directly or indirectly, helps others to develop independently and
distinctively. Herein consists the true justice, which Westermarck rightly calls
``the flower of all moral feelings.''\footnote{``Justice, the flower of all moral
  feelings.'' E.\ Westermarck, \emph{The Origin and Development of the Moral
  Ideas}, 2 vols., London \& New York: Macmillan, 1906--08, vol.~1, p.~124. Cf.\
  my \emph{Etik. En Fremstilling af de etiske Principer og deres Anvendelse paa
  de vigtigste Livsforhold}, Copenhagen: P.G.\ Philipsen, 1887 (and later), III,9;
  IX,1; X,4.} It harmoniously unites self-assertion and devotion, the interest of
the individual and of society. By it society lives in the individuals, not merely
outside and above them. In it the principle of personality comes into its full
right.


% ============================================================
%  V. CAN THE PRINCIPLE OF PERSONALITY BE PROVED
%     AND CARRIED THROUGH?
% ============================================================

\section{Can the Principle of Personality Be Proved and Carried Through?}
\label{sec:5}
% pp. 32--36

\noindent The question can no longer be postponed, whether the principle of
personality, with which we have here occupied ourselves, can be shown to be
valid. In the immediately preceding lecture a justification was indeed given, but
the question is whether this justification does not itself rest on certain
presuppositions. It consisted in this: that the said principle made possible the
most harmonious development of society and the most intimate reciprocal action
between individuals.\footnote{I first gave this demonstration in my \emph{Etik.
  En Fremstilling af de etiske Principer og deres Anvendelse paa de vigtigste
  Livsforhold}, Copenhagen: P.G.\ Philipsen, 1887 (and later), VIII,6, cf.\
  XIII,6, where I also criticized Kant's artificial argumentation. Geijer's
  demonstration (\emph{Samlade skrifter}, sec.~I, vol.~5, Stockholm: P.A.\
  Norstedt \& Söner, 1852, pp.~132--37) rests chiefly on the fact that the human
  being discovers in every other human being a being of the same worth as
  itself, as well as on the fact that the human being can develop into a
  personality only in reciprocal action with other human beings. But this does
  not, after all, exclude that the human being can make these kindred beings,
  whose co-operation it cannot do without, into mere means for itself! --- Karl
  Marx (\emph{Das Kapital. Kritik der politischen Oekonomie}, vol.~1: \emph{Der
  Produktionsprocess des Kapitals}, Hamburg: Otto Meisner, 1867, p.~646) and
  Léon Bourgeois (in \emph{La Philosophie de la solidarité}, Paris: Colin, 1902,
  p.~38) set up the proposition as self-evident, which a glance at history easily
  shows to be incorrect.} But this train of thought is naturally taken as valid
only by those for whom such a development of the human race stands as a good. It
follows from the thought of universal human welfare. But the question is how
matters stand with this thought itself. The answer to this question becomes
decisive for the position of ethics as a science.

Value can be measured only by value. And value points back to a feeling---and
about feelings, it is said, there can be no dispute. Yet one can, when a certain
value is known and acknowledged, infer from it what place other values
consistently occupy in relation to the given value. In every human being---and on
the whole in every people and in every age---a certain fundamental value asserts
itself, an interest that lies deepest and exercises the strongest influence. It
corresponds to, or rather belongs to, what I call the real side of the
personality. And it decides what else acquires value, and what value it acquires.
He for whom his own advantage is the fundamental value cannot, without
inconsistency, sacrifice himself for other persons or for general ends.
Conversely, he in whom love of humanity determines the fundamental value cannot
consistently act egoistically. Our ethical judgments are determined by our
fundamental values, and these in turn by our ruling feelings.

But can one then assert that all human beings, at all times and under all
circumstances, have the same fundamental values? That would surely be an
unpsychological and unhistorical assertion, notwithstanding that thinkers like
Kant and Bentham have, each in his own way, maintained it. The fundamental values
arise through psychological and historical processes that are never present as
wholly completed. A fundamental value can never be set up as self-evident. It
must grow forth, and it depends in each individual case how far the growth has
advanced. It is of no help to say that the human being grows with his greater
ends, for the human being must have grown in order to be able to set himself
greater ends, and must then have grown further by working at them. It is of no
use, as Goethe has said, to speak to the larva that crawls in the dust about how
it shall one day fly from flower to flower.

The principle of personality as a general ethical principle presupposes the
development of a consciousness of humanity, a general love of humankind, such as
appeared in the last centuries before Christ through a peculiar fusion of
Platonic and Stoic ideas. Later it has had a long and changing history. It has
had to struggle, not merely with the egoism of individuals, but also with the
egoism of families, classes, nations, and confessions. In the eighteenth century
it broke fully through and made its definite demands on bourgeois society. But it
is constantly a church militant. Therefore the consequences that are drawn from
the general consciousness of humanity cannot become valid for all either.

Our thought is able to do nothing other than draw the consequence of certain
definite presuppositions. It is the affair of experience to decide whether these
presuppositions are found everywhere and in everyone. This does not diminish the
significance of thought, nor of ethical thinking. Thought can only bring to clear
consciousness what lies in our presuppositions. And the significance of
philosophical ethics rests precisely on this too. It is to set forth the
consequences that can rightly be drawn from a certain stage of the development of
human feeling. Later it can be investigated how far the consequence of other
stages of development coincides with them, and what difference unequal
motivations and unequal presuppositions bring about.\footnote{I developed this
  conception of the scientific limitation and significance of ethics long ago in
  my \emph{Etik. En Fremstilling af de etiske Principer og deres Anvendelse paa
  de vigtigste Livsforhold}, Copenhagen: P.G.\ Philipsen, 1887 (and later),
  chapter III. It is a pleasure to me to meet a similar thought in E.\
  Westermarck's \emph{The Origin and Development of the Moral Ideas}, vol.~1,
  London \& New York: Macmillan, 1906, pp.~17, 104\,ff.}

Where the ethical motivation stands at its limit, because its presuppositions are
no longer acknowledged, there arises the question of how these presuppositions
can be realized. This becomes the affair of moral pedagogy. It is through
education that the fundamental values are formed. The word education is taken here
in the widest sense, so that the whole history of mankind is seen as one great
education. It is the affair of psychology and sociology to show how this
education proceeds involuntarily and unconsciously in individuals and in
generations, and what means stand at our disposal when we will, fully
consciously, to educate human beings.

Schopenhauer has said that it is hard to ground morality, but easy to preach
morality. In the first part of his sentence he is right. But he is wrong in
saying that it should be easy to preach morality, if by this is meant leading
human beings to such a development of their personality that they can, by their
own power and desire, find the right. Neither under education may a personal
being be regarded and treated merely as a means---not even as a means to its own
future. The larva ought to be treated as a larva, right up until it becomes a
butterfly. It is the great thought in Rousseau's pedagogy that the child is to be
treated as a child and not merely as a becoming adult. It is thus a matter of
taking the one who is to be educated at the standpoint at which he stands, and
then, along the paths of self-development and self-activity, leading him to the
level, the fundamental value, that can form the starting point for the inferences
one wishes him to draw. Here Spinoza's words hold, that there is no way in which
one can better show one's spiritual power and capacity than by leading others to
live according to their own clear conviction.

If difficulties are thus bound up with the proof of ethical principles, there are
also difficulties in their being carried through in each individual case. Even if
the fundamental value is presupposed to be the same, the question nevertheless
arises whether, proceeding from the common presuppositions, one can make the same
ethical demands of all individuals. Individuals are different, both with respect
to the kind of their abilities and with respect to the degree of psychic energy
they have at their disposal. From this it follows that one cannot demand the same
of all. Neither what is to be demanded, nor how much is to be demanded, can be
established once and for all, but must be determined with respect to the nature
and powers of the individuals. The same burden weighs differently on different
shoulders. There are human beings who are forced to fight a hard battle to
fulfill the most elementary demands with respect to self-control, which others
fulfill completely involuntarily. The ethical work can be the same, although the
results become of highly different value. Human beings are like runners who begin
the race from very different starting points. For one cannot deny that in the
ethical respect, as in other respects, human beings are very differently equipped
from the beginning.

It would conflict with the principle of personality to judge the individual by an
abstract rule that takes no account of his individual conditions. The task becomes
to individualize the ethical demands, so that each one comes to perform the work
that accords with his personal conditions, and is thereby led onto a track that
can lead further. Ethics is still far from such an individualization, and this
also requires, after all, great psychological and great practical knowledge of
human beings, as well as great patience. It is easier to systematize and to
enunciate universally valid propositions than to find the individual starting
points in the particular cases. System and variation stand, in ethics as in
biology, in sharp opposition to one another.

In concepts such as equity and pardon, however, individualization has long since
been acknowledged. Aristotle already set up the concept of equity as a peculiar
form of justice, necessitated by the fact that there are individual nuances that
prevent the application of a general law. And that pardon is applied can,
juridically, be motivated only by the fact that legislation is not in a position
to take account of all individual circumstances. Pardon is, as Ihering says,
justice's self-correction.

It is, however, not merely downward that individualization is demanded, but also
upward. There are not merely individuals who begin their ethical development from
a lower level than the average; there are also individuals who from the beginning
are equipped with special conditions for ethical work. And by virtue of the
principle of personality the ethical law must make greater demands on them than on
others. They must themselves demand more of themselves than is demanded of
others. They debase themselves by being content with a minimum. It is their
``duty'' to acquire ``merit.'' Schiller's Wilhelm Tell recognizes this when he
answers his wife's complaint---that his countrymen always turn to him when
something great is at stake---with the words: ``Ein Jeder wird besteuert nach
Vermögen!'' [Each is taxed according to his means!]

In the pedagogical respect the distinction between ``duty'' and ``merit,'' on
which Catholic theology lays such great weight, can be of significance. Partly it
can be hard enough to get all to fulfill the minimal demands that one often thinks
of with the word ``duty''; partly it can be needful to warn against ethical
foolhardiness, against venturing further than one's powers allow. The difficulty
lies in this: that only by attempting, by venturing, can we come to know the
extent of our powers. By holding ourselves back we can cause powers and
possibilities in us, which could have been developed through work for greater
ends, to remain unused. Perhaps the attempt to perform a great action leads to our
coming to know our limitation. Perhaps it leads directly to a tragic result. But
such, after all, are the conditions of life. As Longfellow has said: ``Only those
who brave its dangers comprehend its mystery.''


% ============================================================
%  VI. THE PRINCIPLE OF PERSONALITY AND THE VIEW OF LIFE
% ============================================================

\section{The Principle of Personality and the View of Life}
\label{sec:6}
% pp. 37--41

\noindent There still remains to be investigated the significance of the
principle of personality for the view of life and for faith---that is, its
religious significance, when one takes the word ``religious'' in the wider sense.
I believe that this significance is in constant growth, and that the most
essential side of the position of the religious problem in our days rests upon it.

There are three elements that at any given time determine human beings'
faith: science, tradition, and personal experience. They play a very different
role in different individuals, and even where they all assert themselves strongly,
they are combined with one another in very different ways, whereby the view of
life acquires a different timbre. There are scarcely two human beings in whom they
appear in the same manner or are connected in the same way. It would be an
interesting study to penetrate into these inequalities. Unfortunately such a study
is possible only in an imperfect degree with respect to contemporaries.

More and more the first two elements, science and tradition, lead over to the
third, a personal experience of the course and the conditions of life.

(\emph{a})~The newer science, whose origin is to be sought in the century of the
Renaissance, has exercised great influence on the view of life, especially
through two principal thoughts. The one is the conception of the unsurveyability
of the universe, of the unbounded extension of space and time. The bounded stage
on which the human being had hitherto been able to picture the epic of existence
spread out now appeared as a vanishing point. The other is the conception of a
great law-boundness in the whole of nature, a law-boundness that is proved step
by step by the fact that every change in the world of nature or of mind is
explained by its definite connection with other changes, without causes that lie
outside a possible experience being invoked. This scientific concept of cause is
first applied to the inorganic domain, later also to the organic domain, as well
as to the spiritual domain.

The impressions that this great change in the conception of the world made on
human beings' view of life were very different. In some it awakened enthusiasm, in
that existence now, more than before, acquired the stamp of the sublime, at the
same time as the conviction was clarified that it is, after all, the same forces
that work in the nature of the human being and in the infinite universe. In
spirits such as Bruno, Böhme, and Spinoza this direction in the view of life
asserts itself. In others this expansion of the horizon, which transformed the
earthly stage into a vanishing point, awakened anxiety and dread, and led to an
all the more passionate clinging to the old faith. Pascal is the great example of
this direction. There were, finally, those who laid no weight at all on the great
change in the world-picture, who did not become clear about its fundamental
significance, whether in their view of life they held to the traditional faith or
went their own ways. --- Even to this day one finds these three standpoints with
respect to the great problem.

One need not exaggerate the significance of science to be sure that it will come
to exercise an increasing influence on the views of life. One will not be able to
go around science. Pascal already said, after all, that if the Earth really
rotates, then it rotates with both believers and unbelievers. Science will work
most through the quiet power that lies in its method, its way of asking and of
answering. Less and less will human beings who are educated with this method be
able to take mystical and anthropomorphic explanations as real answers to the
riddles of experience. It is less the special results of science than the way of
thinking that it teaches human beings that will come to determine the future of
the views of life. In the first section of this presentation I sought, moreover,
to emphasize that it is no alien power that here comes to meet the personality.
The scientific interest and the scientific capacity are closely connected with the
essence of the personality and bear its stamp. The deep longing for wholeness,
which is the innermost in the personality, also lies at the basis of science.
Science is the work of the personality, which it cannot disown. And it is of no
help to close one's eyes to the expanded horizon that science opens. We must see
ourselves and the stage for all our work and all our interests as smaller circles
within the great circle whose center is everywhere and whose circumference is
nowhere.

Meanwhile it will scarcely ever become possible, from the standpoints and results
of science, to construct views of life for the various individual personalities.
As we have seen (at the end of the second and third sections), it is indeed the
last and highest task of science to understand and explain individual totalities,
and therefore personalities too; but here there can be only approximate
assumptions, and science here approaches art. For the individual, the purely
individual experiences of life will always be the essential thing for the
rounding-off that he involuntarily makes when the view of life is formed. A human
being is more than an example of a general law, and even the human life that moves
in the narrowest circle can give its contribution to mankind's great common
experience. At any rate, only the individual can know how the course of life and
fate is at the particular place in existence where his existence finds itself.
``The humblest defends its place in life's garland,'' a Danish poet has said.

(\emph{b})~In whatever manner the religious traditions may have arisen, their
continued subsistence and the significance they still have for so many rest on the
circumstance that in dogma and cult the deepest experiences of life of countless
generations are deposited. The history of religion shows us how conceptions that
were earlier religiously indifferent have, under the influence of great movements
within the life of feeling, been taken into the service of religion. Thus it has
gone, for example, with the idea of the transmigration of souls in India and with
the thoughts of Greek philosophy in the dogmatics of the Christian church. But
within the religions they are regarded not merely as symbols of experiences of
life. They are fixed as dogmas that are to contain the solution to the riddle of
life. They are torn up from the living soil in which they have grown. Yet more and
more, in the religious consciousness of our time, a conflict asserts itself
between the symbolic-poetic and the literal-dogmatic conception. It comes ever
more clearly forward that the riddles which science cannot solve are not solved by
the religious conceptions either, already for the reason that science both asks
and answers differently than religion does. The time is past when religion could
satisfy all sides of the human life of the soul. Intellectually it has lost the
battle, however matters may stand with the other sides of the life of the soul. In
our days much noble energy is wasted, in the most independent spirits, on the work
that the conflict between religion and science calls forth---an energy that could
have been applied to founding a view of life grounded on one's own experience.

And yet it is not the relation between science and religion that is the most
important thing in the position of the religious problem in our days. From both
sides what can be said has been said, and continued discussion will bring nothing
new. The decisive point lies, by contrast, in the relation between religion and
personality. In opposition to the immutability of the religious tradition stands
the free, the personal experience of life, whose richness and fullness grow as the
distinctive character and independent justification of the individual
personalities are acknowledged. Every human individual is one of the centers where
the value of life is descried, and one of the starting points from which new
values can be found or produced. It is not so much before the tribunal of science
as before that of life that religion is to be placed in the future. This happens by
virtue of the principle of personality. For just as little is it to be a mere
means for the preservation of a religious tradition.\footnote{Since I have in the
  foregoing often referred to Erik Gustaf Geijer as one of the most important
  representatives of the principle of personality, I shall adduce a statement
  showing that he remained true to his principle in his religious conception too.
  ``I am,'' he writes in a letter, ``neither a church-Christian nor even a
  Bible-Christian \ldots\ [I] am a Christian on my own account. Thousands find
  themselves in the same position as I---in fact all who cannot content themselves
  with a mere authority-faith, and who are compelled to build their own religion
  out of their own experience and reflection. This individual element, now the
  only living thing in religion \ldots\ points everywhere, however, toward a
  common goal [\ldots].'' (\emph{Samlade Skrifter}, sec.~I, vol.~1, Stockholm:
  P.A.\ Norstedt \& Söner, 1849, p.~XXXII\,f.) (Emphasis mine.) Yet Geijer, who
  stood under the strong influence of the romantic speculation of his time, also
  gave his principle of personality a theological and metaphysical significance
  which is untenable, since every attempt in this direction leads to a play with
  analogies and allegories that can neither be carried through consistently nor
  confirmed by experience. Cf.\ on this my \emph{Religionsfilosofi}, Copenhagen
  \& Kristiania: Gyldendalske Boghandel, Nordisk Forlag, 1901 (and later),
  \S32.}

The religious need is a need to hold fast the values of life beyond the limits
within which the human will can work for them. Could the human being himself
completely secure his spiritual and material property absolutely, religion would
not arise in any form. But now there comes hope or fear, exaltation or dejection,
according as one feels oneself at the high point of life or in its depths. The
cults of the great religions give, in the most exalted poetry, expression to these
opposite moods that the epic of life brings with it, and which will always assert
themselves in everyone who lives life wholly and fully, and whose sympathy is not
blunted and limited.

I will not prophesy about the religion of the future. I am a philosopher, not a
prophet. But everything speaks for it, that the more the real and constant
experiences of life acquire influence on the religious life, the more the
symbolic-poetic conception of the religious experiences will replace the
literal-dogmatic one, and make it ever more conscious. And thereby a great step
will have been taken toward a free, self-acquired human view of life. ---

(\emph{c})~Both science and tradition thus point to personal experience of life. A
thoroughgoing individualization of the view of life is necessary. The personality
is to set its stamp on its faith. Only thus does this faith become a testimony to
how life moves at this particular place in the world. And when advanced
psychological investigations lead us to recognize that the individual differences
are far greater and deeper than hitherto assumed, the views of life of the future
will probably exhibit greater oppositions than one can get a glimpse of so long as
the common confessional dogmas cover the mutually divergent conceptions and
applications of an identically worded confession.\footnote{Cf.\ the survey, of the
  individual differences that can acquire significance for the view of life, which
  I have sought to give in my \emph{Religionsfilosofi}, Copenhagen \& Kristiania:
  Gyldendalske Boghandel, Nordisk Forlag, 1901 (and later), \S\S36--43 and
  \S\S94--96. --- We stand here at a point that has been strongly emphasized by
  Danish philosophers, all the way from Holberg to Sibbern and Kierkegaard. See on
  this my book \emph{Danske Filosofer} [Danish Philosophers], Copenhagen:
  Gyldendalske Boghandel, Nordisk Forlag, 1909.} It is a gleam of light in the
darkness of the time that this side of the principle of personality is more and
more being acknowledged. If it is acknowledged completely, a wholly new conception
of the deepest problems will assert itself. For the time being the old dogmas
struggle with the new doubt. Through this purgatory the decisive victory is to be
won for the principle of personality.

\end{document}
